\section{A distribution where the important coordinates are known}
\label{sec:controlled}
Real embeddings did not establish an advantage for B or C in the independent comparison.
To understand the mechanism, we can construct a distribution in which the score tells us
exactly which coordinates matter most. This is an experiment about the algorithms, not a
claim that a trained Matryoshka embedding has this distribution.

\subsection{State the assumptions before calculating recall}
Let document signs be independent and balanced: each coordinate is independently $+1$ or
$-1$ with probability $1/2$. A query has $m$ strong coordinates of magnitude $a$, and
$d-m$ weak coordinates of magnitude $\epsilon>0$. Its signs are independently random.
Scaling the whole query to unit length does not change any ranking or the following inequality:
\begin{equation}
(d-m)\epsilon<a.
\label{eq:strong-dominance}
\end{equation}
A single strong mismatch loses $2a$ score. Even changing every weak coordinate from its best
to its worst sign loses only $2(d-m)\epsilon$. Therefore any document matching all strong
coordinates outranks every document with a strong mismatch. More generally, fewer strong
mismatches always outrank more strong mismatches; the weak coordinates decide the order inside
one mismatch count.

The experiment sets $d=256$, $m=13$, $a=1$ and $\epsilon=0.0001$. The weak total is $0.0243$,
so Equation~\eqref{eq:strong-dominance} holds. One setup puts the strong coordinates first.
The other chooses a different set of 13 coordinates for each query. Neither changes the document
distribution. All 40 generated queries are retained, with the first 20 used for tuning and the
remaining 20 for checking frozen choices.

\subsection{Recall from opening only the ideal strong key}
Let $Z$ count documents matching all $m$ strong signs. Under the stated independent distribution,
\[
p=2^{-m},\qquad Z\sim\operatorname{Binomial}(N,p),\qquad \E[Z]=N2^{-m}.
\]
Suppose a method opens exactly that key, scores every document in it, and returns at most $K$.
Every one of those documents outranks all documents outside the key. Its recall on this corpus
is therefore exactly
\begin{equation}
R_{\rm ideal}=\frac{\min(Z,K)}{K}.
\label{eq:ideal-key-recall}
\end{equation}
For the expected recall over random corpora, use the fact that $\min(Z,K)$ counts one for every
threshold $i\in\{1,\ldots,K\}$ that $Z$ reaches. Taking expectations gives
\begin{equation}
\E[R_{\rm ideal}]=\frac{1}{K}\sum_{i=1}^{K}\Pr(Z\ge i).
\label{eq:expected-ideal-recall}
\end{equation}
This is a finite binomial calculation. It does not replace $Z$ by its average and then assume
that every query has that many matching documents.

\begin{examplebox}{Why a million documents can make a short key more useful}
At $m=13$, a 100,000-document corpus has about 12.2 ideal matches on average. Opening only
that key cannot usually return the top 100. At one million documents, the average is about
122.1, so one key often contains enough results.

In the independent fixed-coordinate sample, 19 queries had at least 100 ideal matches and
one had 99. Equation~\eqref{eq:ideal-key-recall} predicts mean recall
$(19\times100+99)/(20\times100)=99.95\%$. The scan that opened only that cluster measured exactly
99.95\%. No fitted embedding-accuracy percentage was needed.
\end{examplebox}
If the strong coordinates move between queries, a fixed first-$m$ key may contain only a few
of them. Equation~\eqref{eq:ideal-key-recall} then does not describe that key's recall. The
assumption is about the coordinates actually indexed, not simply the number of indexed bits.

\subsection{How many keys must exact lookup try?}
For a key containing exactly the $m$ strong coordinates, equal strong magnitudes make all keys
with $r$ mismatches share a prefix penalty. Let $Z_r$ be the number of documents whose strong
key has exactly $r$ mismatches. For a fixed query,
\[
\Pr(\text{one document has }r\text{ strong mismatches})=\binom mr2^{-m}.
\]
The vector of counts follows a multinomial distribution; the counts at different radii are
not independent. Let $r_*$ be the first radius with at least $K$ documents in total:
\[
r_* = \min\left\{r:\sum_{i=0}^{r}Z_i\ge K\right\}.
\]
The measured exact key search finishes the tied key layer at that radius. Under
Equation~\eqref{eq:strong-dominance}, it can then reject all larger radii. Its work is
\begin{equation}
H=\sum_{i=0}^{r_*}\binom mi,\qquad F=\sum_{i=0}^{r_*}Z_i.
\label{eq:strong-key-work}
\end{equation}
$H$ includes empty keys. $F$ counts full document scores in the present implementation.
For $r_*=1$ and $m=13$, that means 14 key attempts, not one attempt and not $2^{13}$ attempts.
A smaller lookup limit invalidates the exact-stop conclusion. A candidate quota can also stop
before the layer is complete, so its recall must be measured using the actual stopping rule.

The probability that radius $r$ contains enough documents is another binomial tail:
\[
\Pr(r_*\le r)=\Pr\!\left(\operatorname{Binomial}\!\left(N,\,2^{-m}\sum_{i=0}^{r}\binom mi\right)\ge K\right).
\]
This supplies a distribution over work, rather than declaring that the expected occupancy is
an exact count for every query. No independence assumption between the different radii is needed.

\subsection{Choose how far bitplanes should split}
B orders the coordinates using the actual query magnitudes, so it can follow the strong signs
even when their positions change. Suppose at least $K$ documents match all $m$ strong signs.
After following $j\le m$ preferred splits, those $K$ correct results are still present. Score the
remaining set and the full-score bound can reject the alternatives that missed a strong bit.

With $W=\ceil{N/w}$ physical words per mask, the one-path approximation is
\begin{equation}
T_B(j)\approx T_{\rm prepare}+jWc_{\rm split}+Wc_{\rm leaf}
                    +N2^{-j}c_{\rm score}+T_{\rm queue}(j).
\label{eq:one-path-time}
\end{equation}
Both split and leaf loops keep the original width $W$. For an actual corpus, replace $N2^{-j}$
by the observed surviving count. If fewer than $K$ documents match all strong coordinates, more branches may
be required and this single-path expression is insufficient; those queries stay in the study.

One more cut saves about $N2^{-(j+1)}$ document scores. It helps only when those saved scores
cost more than another full-width split and the added queue work:
\begin{equation}
N2^{-(j+1)}c_{\rm score}>
 Wc_{\rm split}+T_{\rm queue}(j+1)-T_{\rm queue}(j).
\label{eq:one-more-cut}
\end{equation}
Ignoring the smaller queue term gives a first estimate from the derivative:
\[
\frac{dT_B}{dj}=Wc_{\rm split}-N\ln(2)\,2^{-j}c_{\rm score}=0,
\qquad
j_c=\log_2\!\left(\frac{Nc_{\rm score}\ln 2}{Wc_{\rm split}}\right).
\]
The second derivative of the score term is positive. Check the neighboring integers and
boundaries $0$ and $m$, then check real leaf sizes near $N2^{-j}$. This optimizes one stated
model. It does not optimize arbitrary clustering, discrete query paths and hardware behavior
with one derivative.

When $N$ is large, $N/W$ is about the fixed word width $w$. Thus $j_c$ need not grow with $N$.
A fixed leaf size can be a poor scaling choice: it forces extra cuts as the corpus grows,
even when another full-width mask pass costs more than the document scores it removes.
Both global dense-bitplane traversal and fixed-dimension scan still have at least linear
work in $N$ in this model. B can improve the constant; the first full-width pass does not
become a logarithmic document lookup.

For $j\ge1$ with nonempty alternatives, the last split can hold $j+2$ full masks: queued
alternatives, the parent and both children. The corresponding mask payload is
\[
M_{\rm masks,peak}=(j+2)Wb_w.
\]
With no split, only the root mask is needed. Queue metadata, score tables, outputs and the
persistent index are additional memory.

\subsection{What happened when the strong coordinates stayed fixed}
All methods could use direct-key routing as well as the centroid and sign-prefix scans.
There were 704 tuning configurations and 27 independent evaluation workers. The fastest tested
key configuration used one physical cluster and a 13-bit key. The scan baseline chose direct
routing into 8,192 clusters, which uses the same 13 strong signs.

At one million documents, C measured 0.0127 ms median latency and 100\% binary recall. A
measured 0.0145 ms and 99.95\% recall. Their paired speedup estimate was 1.145, with a 95\%
interval of about 1.095 to 1.196. Both met every requested target through 99\%.
At 100,000 documents and the 99\% target, both reached 100\% recall: C measured 0.0151 ms
and A 0.0186 ms, a paired estimate of 1.227 with interval about 1.144 to 1.364.

The difference is a few microseconds, and the methods largely overlap in candidate selection.
C's key enumeration and scoring share one native call. A first ranks cluster keys through a
native call, then scans their rows through another. The measurements establish an implementation
advantage for this integrated lookup, not a theorem that key lookup dominates a scan baseline
that incorporates the same candidate-selection rule. Larger cluster and key grids may also
change the fastest tested setting.

\subsection{What changed when the strong coordinates moved}
The second setup uses the same document signs, but chooses each query's strong coordinates
uniformly without replacement. Fixed-prefix keys cannot follow those changing positions.
B can, because its coordinate order comes from the query itself.

This run measured 1,001 tuning settings and 60 independent evaluation workers. At one million
documents and the 99\% requirement, the frozen B configuration used one cluster, unlimited
traversal and leaf size 128. It reached 100\% binary recall at 0.1496 ms median latency.
The scan configuration used 256 centroid clusters and 255 probes, reaching 99.95\% at
24.9377 ms. Their paired speedup estimate was about 166.7, with a 95\% interval of
147.6 to 183.9. C reached 100\% at 26.8429 ms using one cluster and a 12-bit key.

At 100,000 documents, the 99\% target comparison reached 100\% recall for all three:
A measured 2.5438 ms, B 0.1375 ms, and C 2.7170 ms. B's paired speedup over A was about
18.5, with a 95\% interval of 18.1 to 19.0. These are selected-settings comparisons within
the declared grid, with query and process-block uncertainty. They are not proofs of a global
optimum over every cluster count, leaf size, routing algorithm or implementation.

\fig{measured-cases}{Independent measurements at one million documents, a 99\% binary-recall requirement and a 32 GB memory limit. Labels show median time and achieved recall. The vertical axis is logarithmic: adjacent major ticks differ by a factor of ten. The real-data evaluation has 200 queries; each constructed case has 20.}{1}

\subsection{Compare predicted work with the saved native counters}
The count checker uses document signs recorded during preparation, independently of native
index traversal. It checks the B's single-path formula only when at least $K$ documents match all strong
coordinates and a leaf is reached before splitting weak coordinates. For C, it checks the
single-key bound case and Equation~\eqref{eq:strong-key-work} when the key contains all strong
coordinates. Empty keys are counted, including keys in the final tied layer.

Across the completed tuning and evaluation runs, 447 applicable per-query predictions matched
the native counters exactly. This includes split words, leaf words, full scores, key attempts
and the tracked B mask peak where applicable. Other queries and clustered configurations remain
in the benchmark; the checker reports them as outside these specific prediction conditions.
It does not silently treat them as successful predictions.

There are two different results here. The component and conditional-work checks support the
stated formulas. The timed experiment establishes how the implemented methods performed on
these data. Neither alone establishes a universal latency forecast. In particular, if a new
machine changes the relative cost of a bitmap pass and a full score, Equation~\eqref{eq:one-more-cut}
can favor a different leaf size.

The complete settings and raw evidence live under
\path{results/controlled-fixed-2026-09-10/} and
\path{results/controlled-adaptive-2026-09-10/}. The small runnable example is
\path{examples/small_search.py}; it prints one query, its stored document bits, exact results
and each method's work. \path{experiments/check_controlled.py} connects the conditional formulas
to the recorded counters.
